\documentclass{drpaper}

\title{Introduction: The Great AI Experiment}
\author{Daniel Rothschild}
\date{\today}

\begin{document}

\maketitle

\section{The great AI experiment}

Recent advances in artificial intelligence constitute one of the most remarkable and consequential scientific developments of our era. In a few short years, computers have acquired the ability to write original prose, engage in sustained conversation, recognise objects in photographs, generate images from text descriptions, and defeat human grandmasters at games of extraordinary complexity. These are not incremental improvements on older capacities. They represent a qualitative shift in what machines can do — a shift that has surprised even those who built the systems responsible for it.

This book takes these developments seriously as scientific evidence. The construction and deployment of large-scale AI systems is, I want to argue, an experiment — perhaps the most ambitious experiment ever conducted — about the nature of learning. By building machines capable of acquiring the kinds of skills that humans acquire through childhood development and education, we are generating indirect but powerful evidence about how learning works. We have, for the first time, systems whose training we control, whose internals we can inspect, and whose performance we can measure with precision. That is an unprecedented vantage point on a question that has puzzled philosophers and scientists for millennia.

The question this book addresses is: what does the great AI experiment tell us about learning — in machines and, by inference, in biological minds?

This is not primarily a book about AI. It is a book about learning, which uses AI as its principal source of evidence. The distinction matters. Many books in this space either explain how AI systems work for a general audience, or extrapolate from AI to speculative claims about superintelligence and existential risk. My aim is different and more modest: to use what we can observe about how AI systems learn as a window into the nature of learning itself. The results are philosophically significant, and they bear on debates in cognitive science, philosophy of mind, and epistemology that long predate the current wave of AI.

\section{What AI tells us about learning}

The book defends four main claims.

\subsection{Learning as search}

The first and most general claim is that learning, across all its forms, is best understood as \emph{search}. To learn is to navigate a space of possible systems — possible ways of representing and responding to the world — guided by experience. This framework is broad enough to encompass Bayesian updating, the gradient-based training of neural networks, and the developmental trajectories of human children. It is also old enough to have been anticipated, in different forms, by Plato and by the founders of modern AI \citep{NewellSimon1976}.

The learning-as-search framework does immediate philosophical work. It dissolves the ancient puzzle that Plato raised in the \emph{Meno}: how can you seek what you don't already know? If learning is search through a pre-existing space of possible systems, guided by a generator and a test, then new knowledge can be extracted from the environment without being presupposed. The puzzle assumed that learning required identifying a specific target in advance; the search framework shows why it doesn't \citep{PlatoMeno}.

The framework also connects biological and artificial learning in a principled way. Both brains and neural networks learn by searching through parameter spaces — vast landscapes of possible configurations — guided by error signals. \citet{dehaene2020we} makes this explicit: the brain's learning mechanisms are, at the most abstract level, a massive search through parameter space driven by prediction error and reward. The same description applies to the training of a deep neural network. This does not mean the mechanisms are identical at a finer level of description. But it locates them in the same conceptual space, which is the first step toward a general theory.

\subsection{Supervised learning as the engine of AI}

The second claim is more specific: \emph{supervised learning} — training on labelled input-output pairs — is the unifying mechanism behind the apparent diversity of modern AI's successes. Language models, image generators, game-playing systems, and protein-folding algorithms all look very different on the surface. What I argue is that, under the hood, they all reduce to the same basic computational operation: adjust parameters so as to reduce prediction error on a training set, and generalise to new inputs.

This claim may seem surprising. Language models appear to learn language, not to perform a pattern-matching task. Image generators appear to be creative. Game-playing systems appear to reason strategically. But in each case, the breakthrough depended on recognising that the problem could be reformulated as supervised learning — finding the right training signal that would allow gradient-based optimisation to do its work \citep{deeplearninglecun}. This reformulation was not trivial engineering; it was a conceptual achievement. And the fact that one mechanism suffices across such wildly different domains is itself a substantive finding about the nature of learning.

\subsection{Modest associationism}

The third claim is philosophical: the success of supervised learning in AI provides support for a form of \emph{associationism} — the view that learning proceeds by the gradual formation and adjustment of associations in response to experience. Associationism has a long history in philosophy and psychology, and an equally long history of being dismissed as too simple to account for sophisticated cognition. The central objection was that rule-based, symbolic processes were needed to explain structured thinking, language, and one-shot generalisation — and that associationist mechanisms simply could not do this work.

The evidence from AI undermines this objection. The same class of mechanisms that critics conceded might handle low-level pattern recognition turns out, when properly implemented and scaled, to handle much more sophisticated cognitive tasks. This is a vindication of the associationist programme, though I argue it is a \emph{modest} one. The successes of deep learning support two core associationist theses: that a \emph{uniform} mechanism operates across diverse learning tasks, and that learning proceeds by \emph{gradual, error-driven adjustment} of something like weighted connections. But the architectures within which supervised learning operates — transformers, convolutional networks, diffusion models — go well beyond anything associationists envisaged. And the behaviour these systems exhibit at inference time looks, in many respects, more like reasoning than conditioning. The upshot is a picture that is associationist about the engine while remaining agnostic about the surrounding machinery.

This also has implications for the empiricism-nativism debate \citep{Buckner2023-BUCFDL}. The domain-generality of supervised learning as a mechanism provides some support for empiricist claims about the uniformity of learning. But it does not settle the question of how much domain-specific structure biological learners bring to learning tasks. These are compatible hypotheses, and the evidence from AI speaks to mechanism, not to whether the mechanism operates on innately structured inputs.

\subsection{Language and thought}

The fourth claim concerns the role of language in cognition. Among all the AI systems developed in recent decades, only those trained extensively on natural language exhibit powerful domain-general reasoning. Systems trained on images, games, or other structured data without natural language show impressive but narrow competence. This is a striking empirical result, and it demands explanation.

My explanation is that natural language, as a system of representation, has properties that make general inference computationally tractable. Language is abstract: it compresses information across domains into a shared representational medium. It is productive: new concepts can be constructed compositionally from existing ones. And it is, crucially, \emph{public}: training on language means training on the accumulated inferential practices of a linguistic community, not just on perceptual statistics. These properties collectively make language a uniquely powerful substrate for general reasoning \citep{lupyan2016language}.

This result has implications for the longstanding debate about the relationship between language and thought. One tradition, associated with Fodor and others, holds that thought is prior to language — that we think in a language of thought, and natural language merely expresses thoughts already formed. Another tradition holds that language actively shapes and transforms cognition, rather than merely expressing it. The evidence from large language models supports the second view. Adding language to an AI system does not merely give it a new output channel. It transforms what the system can do \citep{MAHOWALD2024517}.

\section{Overview of the book}

The book is structured to develop these four claims in sequence, moving from the most general to the most specific.

Chapter 1 develops the learning-as-search framework in detail, tracing its roots in classical philosophy and early AI, and showing how it applies to contemporary machine learning. It also sets up the central philosophical question: what can learning from experience achieve, and where are its limits?

Chapter 2 examines learning as a computational process — asking what the theory of computation tells us about what learning can and cannot accomplish. This includes a discussion of Valiant's PAC learning framework and the limits it identifies, before turning to the question of why deep learning succeeds in ways that PAC theory would not have predicted.

Chapter 3 develops the taxonomy of machine learning paradigms — supervised, unsupervised, and reinforcement learning — and argues that supervised learning is the central mechanism linking them. This chapter also introduces deep learning and the gradient-based training that makes it work.

Chapter 4 turns to reinforcement learning and motivation, asking what the reward signal is for biological learners. Is understanding itself intrinsically rewarding? And can any formal objective function capture the values that actually drive human learning?

Chapter 5 examines the associationist interpretation of deep learning — what its successes vindicate, and where the vindication runs out. This is where the relationship between training and inference becomes central, and where the limits of the associationist picture come into view.

Chapter 6 takes up the language thesis in detail: why language-trained systems are unique, what this tells us about the cognitive utility of language, and what it implies for debates about linguistic relativity and the relationship between language and thought.

Throughout, the aim is to use the evidence from AI to illuminate questions about learning that are of intrinsic philosophical and scientific interest — questions that were live long before the first neural network was trained, and that the great AI experiment has given us new tools to address.

\bibliography{/Users/daniel/Dropbox/Latex\ files/bib\ files/danbib}

\end{document}
